From $E = h f = \frac{h c}{\lambda}$ the wavelength is $\lambda = \frac{h c}{E}$, with $E$ in joules ($\SI{1}{\electronvolt} = \SI{1.602e-19}{\J}$). Visible light runs from about \SI{380}{\nm} (violet) to \SI{750}{\nm} (red).

\begin{abcliste}
\abc
\begin{align}
 \lambda_1 &= \laF \\
   &= \frac{\nch\times\ncc}{\EaO\times\SI{1.602e-19}{\J\per\electronvolt}} \\
   &= \la \approx \result{\laP{-9}{3}}
\end{align}
$\result{\text{Visible, green}}$ light.

\abc In electronvolts, $E_2 = \EbvP{0}{3}$. The wavelength:
\begin{align}
 \lambda_2 &= \lbF \\
   &= \frac{\nch\times\ncc}{\Eb} \\
   &= \lb \approx \result{\lbP{-9}{3}}
\end{align}
Shorter than \SI{380}{\nm}: $\result{\text{ultraviolet}}$.

\abc
\begin{align}
 \lambda_3 &= \lcF \\
   &= \frac{\nch\times\ncc}{\EcO\times\SI{1.602e-19}{\J\per\electronvolt}} \\
   &= \lc \approx \result{\lcP{-9}{3}}
\end{align}
Longer than \SI{750}{\nm}: $\result{\text{infrared}}$.
\end{abcliste}
The larger the photon energy, the shorter the wavelength. A shortcut: $h c = \SI{1240}{\electronvolt\nm}$, so $\lambda = \SI{1240}{\electronvolt\nm}/E$.
